% ═══════════════════════════════════════════════════════════════════
% ЧАСТЬ XI. АЛГОРИТМЫ ПРОТОКОЛА V1–V9 И ТАБЛИЦЫ ПЕРИОДОВ
% ═══════════════════════════════════════════════════════════════════
\part{Алгоритмы протокола V1--V9 и таблицы периодов}

\section{Пошаговые алгоритмы проверок}\label{sec:algoV}

Здесь описан алгоритм каждой проверки протокола на уровне,
достаточном для независимой реализации (собственно реализация ---
\texttt{laboratory.py}, пункты меню «Полный протокол» и «Стенд
N=15/30»).

\subsection{V1: перепись характеров (точная)}

\begin{enumerate}[leftmargin=2em, itemsep=0.1em]
\item Для $N\in\{15,30\}$ перебрать все пары $(a,b)$ с
  $1\le a$, $1\le b$, $a+b\le N-1$.
\item Вычислить $g_0=\gcd(N,a,b)$ и проводник $d=N/g_0$.
\item Увеличить счётчик $h_d$.
\item Сверить суммы $\sum_d h_d$ с родами $91$ и $406$.
\item Сверить таблицу $h_d$ с формулой Мёбиуса~\eqref{eq:mobius}.
\end{enumerate}
Сложность $O(N^2)$; всё в целых числах; независимость --- две схемы
(перебор и Мёбиус) и целочисленное совпадение с родом.

\subsection{V2/V3: замкнутая форма против tanh-sinh}

\begin{enumerate}[leftmargin=2em, itemsep=0.1em]
\item Построить детерминированную выборку характеров: до трёх на
  проводник (первый, средний, последний), по три данных обхода
  $(r,s)$ на каждый.
\item Для каждой тройки вычислить замкнутую форму
  $P^{\Gamma}=\frac1N\zeta^{ra+sb}\Omega_{a,b}$ при $dps=35$.
\item Независимый эталон: разрезать бета-интеграл заменой
  $t=u^N/2$ на двух половинах; проинтегрировать tanh-sinh
  (без $\Ga$-функций); умножить на константу обхода
  $\frac1N e^{2\pi i(ra+s(b-N)\bmod N)/N}$.
\item Относительное отклонение $|P^\Gamma-P^{tanh}|/\max(|P|)$
  записать в таблицу по проводникам; фазу $\arg P^{tanh}$ сравнить
  с $2\pi(ra+sb)/N$ по модулю $2\pi$.
\end{enumerate}
Порог PASS: отклонение $<10^{-30}$; фазовое отклонение $<10^{-25}$.
Фактические максимумы: $3.9\cdot10^{-36}$ и $0.0$.

\subsection{V4: эквивариантность}

\begin{enumerate}[leftmargin=2em, itemsep=0.1em]
\item Для каждого характера выборки взять базовый обход $(r,s)$ и
  сдвиги $(u,v)\in\{0,1\}^2$.
\item Сравнить $P(r+u,s+v)$ с $\zeta^{ua+vb}P(r,s)$ --- и в
  замкнутой форме, и численно (два источника).
\end{enumerate}
Фактический максимум $3.9\cdot10^{-36}$; тождество
следствие~\ref{cor:equivar} доказано.

\subsection{V5: целостность цепей}

\begin{enumerate}[leftmargin=2em, itemsep=0.1em]
\item Взять полные степени $h=x^2y^3$ и $h=x^3y$ --- рациональные
  функции с голоморфными дифференциалами $dh$ на кривой.
\item Интеграл $\int_\gamma dh$ по любому замкнутому циклу равен
  $h(\text{конец})-h(\text{начало})=0$.
\item Численно проинтегрировать $dh$ по циклам выборки;
  максимумы модуля записать.
\end{enumerate}
Фактический максимум $1.5\cdot10^{-36}$; проверка ловит ошибки
ветвления при движении по циклам.

\subsection{V6: лестница отражения}

\begin{enumerate}[leftmargin=2em, itemsep=0.1em]
\item Для всех $k=1,\dots,N-1$ вычислить
  $\Ga(k/N)\Ga(1-k/N)$ и $\pi/\sin(\pi k/N)$.
\item Сравнить относительные отклонения по всем $k$ обоих уровней.
\end{enumerate}
Фактический максимум $7.0\cdot10^{-36}$; тождество --- лемма
\ref{lem:reflection} с полным выводом в приложении~\ref{app:gamma}.

\subsection{V7: полнота системы}

\begin{enumerate}[leftmargin=2em, itemsep=0.1em]
\item Построить матрицу $M$ размера $g\times N^2$:
  $M[(a,b)][(r,s)]=\zeta^{ra+sb}\Omega_{a,b}$ по всем парам.
\item Вычислить численный ранг (SVD с порогом $10^{-25}$) и число
  обусловленности $\sigma_g/\sigma_1$.
\end{enumerate}
Фактически: ранг $91$ и $406$; обусловленности $8.63$ и $18.17$ ---
матрицы хорошо обусловлены, порождающий набор полон без PSLQ.

\subsection{V8: ортогональность ДПФ}

\begin{enumerate}[leftmargin=2em, itemsep=0.1em]
\item Вычислить $\sum_{r,s}P_\chi\overline{P_{\chi'}}$ по всем парам
  характеров выборки.
\item Сравнить с $\delta_{\chi\chi'}|\Omega_\chi|^2$.
\end{enumerate}
Фактически: диагональ $4.0\cdot10^{-35}$, внедиагональ
$8.2\cdot10^{-37}$ --- прямое следствие полноты системы характеров
группы $(\ZZ/N)^2$.

\subsection{V9: глубокая запись}

\begin{enumerate}[leftmargin=2em, itemsep=0.1em]
\item Повторить одну репрезентативную запись при $dps=70$.
\item Сравнить с эталоном; относительное отклонение $1.7\cdot10^{-71}$
  --- согласовано с пределом $dps=70$.
\end{enumerate}

\section{Таблицы значений $\Omega_{a,b}$ и периодов}\label{sec:periodtables}

Таблицы ниже --- численные значения замкнутой формы для
репрезентативных характеров ($dps=35$; точность последнего знака
контролирована V2/V9). Значения $\Omega_{a,b}$ вещественны
положительны; периоды комплексны с точной фазой.

\subsection{$N=15$: характер проводника $d=3$}

Единственный характер: $(a,b)=(1,5)$ --- $a+b=6$;
$\Omega_{1,5}=\Ga(\tfrac1{15})\Ga(\tfrac5{15})/\Ga(\tfrac6{15})$:
\[
  \Omega_{1,5}=15.9787\ldots
\]
Периоды $\frac1{15}\zeta_{15}^{r+5s}\Omega_{1,5}$ при
$(r,s)\in\{(0,0),(1,0),(0,1)\}$:
\[
  P(0,0)=1.06525\,(1+0i);\quad
  P(1,0)=1.06525\,e^{2\pi i/15};\quad
  P(0,1)=1.06525\,e^{10\pi i/15}.
\]

\subsection{$N=15$: характеры проводника $d=5$}

Шесть характеров; первые три: $(1,1)$, $(1,2)$, $(1,3)$.
\begin{center}
\small
\begin{tabular}{@{}l l l@{}}
\toprule
$(a,b)$ & $a+b$ & $\Omega_{a,b}$ \\
\midrule
$(1,1)$ & $2$ & $\Ga(1/15)\Ga(1/15)/\Ga(2/15)=246.897\ldots$ \\
$(1,2)$ & $3$ & $\Ga(1/15)\Ga(2/15)/\Ga(3/15)=36.6265\ldots$ \\
$(1,3)$ & $4$ & $\Ga(1/15)\Ga(3/15)/\Ga(4/15)=11.6686\ldots$ \\
\bottomrule
\end{tabular}
\end{center}

\subsection{$N=15$: характер $a+b=15$ (чистая доля $\pi$)}

Для $(a,b)=(7,8)$: $\Omega_{7,8}=\pi/\sin(7\pi/15)$ ---
чистая дробная доля $\pi$ по лемме~\ref{lem:relations}(1):
\[
  \Omega_{7,8}=\frac{\pi}{\sin(84^\circ)}=3.16298\ldots
\]
Этот блок демонстрирует теорему~\ref{th:normalization}(2) в чистом
виде: трансцендентная часть --- ровно $\pi$ с рациональным углом
$k\pi/N$.

\subsection{$N=30$: характеры проводников $d=6$, $d=10$}

\begin{center}
\small
\begin{tabular}{@{}l l l@{}}
\toprule
Проводник & Пример $(a,b)$ & $\Omega_{a,b}$ \\
\midrule
$d=6$ & $(1,4)$ & $\Ga(1/30)\Ga(4/30)/\Ga(5/30)=198.075\ldots$ \\
$d=10$ & $(1,8)$ & $\Ga(1/30)\Ga(8/30)/\Ga(9/30)=52.9878\ldots$ \\
$d=30$ & $(1,1)$ & $\Ga(1/30)^2/\Ga(2/30)=557.765\ldots$ \\
\bottomrule
\end{tabular}
\end{center}

Значения уровня $30$ с нечётными аргументами (например
$\Ga(1/30)$ в блоке $d=30$) демонстрируют несводимость к уровню
$15$ (теорема~\ref{th:H3}): аргумент $1/30$ не редуцируется ни
отражением, ни рекуррентностью к знаменателю $15$.

\section{Схема полного конвейера: от тройки до индексов}
\label{sec:pipeline}

Сведём весь конвейер программы в одну схему --- это
«карта» монографии в шести строках:
\begin{enumerate}[leftmargin=2em, itemsep=0.15em]
\item Геометрия $\to$ тройка $(n,B,\lambda_0)$: порядок поворота,
  целочисленный индекс, спектральный порог.
\item Тройка $\to$ поправки: $\delta=\pi/n$, $\gamma=\delta^4/B$,
  $\delta_{\mathrm{eff}}=\delta^5/B$, $\Delta_{Ch}$
  по формуле~\eqref{eq:deltaCh}.
\item Поправки $\to$ сертификаты: каждая величина получает
  независимую проверку (точная арифметика / прямое интегрирование).
\item Сертификаты $\to$ решётки: поляризация, SNF-тип, индексы
  блоков --- всё в целых числах.
\item Решётки $\to$ $\Ga$-нормировка: объёмы и модули ---
  $\Ga$-мономиалы, фазы --- в полях классов.
\item $\Ga$-нормировка $\to$ сводный вердикт: протокол V1--V9,
  мультиязычная верификация, ядро Lean.
\end{enumerate}
Каждая строка схемы раскрыта в соответствующей части монографии;
схема демонстрирует замкнутость: из целочисленной геометрии через
алгебру и анализ --- снова к целочисленным инвариантам
(дискриминант $1125^2$, тип $(1,1,5,5,15,15,15,15)$, $Q=480$),
которые и служат финальными сертификатами.
